\documentclass[sigconf,nonacm]{acmart}
\usepackage[UTF8,scheme=plain,fontset=fandol]{ctex}
\usepackage{amsmath,booktabs,array}
\settopmatter{printacmref=false,printccs=false,printfolios=true}
\newcommand{\RRF}{\operatorname{RRF}}
\newcommand{\Rel}{\operatorname{Rel}}
\newcommand{\Recall}{\operatorname{Recall}}
\newcommand{\DCG}{\operatorname{DCG}}
\newcommand{\AP}{\operatorname{AP}}
\begin{document}
\title{排序融合何时有效？\\互惠排名融合的原理、反例与可复现评价设计}
\subtitle{《信息检索原理》第九周论文阅读报告}
\author{Course Project}
\renewcommand{\shortauthors}{Course Project}
\affiliation{\institution{信息检索原理课程}\city{Beijing}\country{China}}
\begin{abstract}
多个检索器给出不同答案时，融合排序既可能利用互补信息，也可能重复放大同一种错误。本文围绕互惠排名融合（Reciprocal Rank Fusion，RRF）的原始论文展开批判性阅读，并结合现代混合检索融合分析与BEIR评价框架，讨论排名、分数、候选集合和评价协议之间的关系。本文将文献证据与自主推导分开：对文献仅概述其可核实结论；通过自行构造的小型排序、参数翻转、分数缩放和重复投票案例，解释RRF的行为与失败条件；进一步给出不依赖测试集调参的课程项目方案。分析表明，RRF的工程吸引力来自其对分数尺度的独立性，而不是对检索器质量、候选深度或参数选择的独立性。合理评价应同时固定输入预算，报告单路基线、相关性、召回上界和查询级差异。本文不将设计中的项目实验作为已完成结果。
\end{abstract}
\keywords{信息检索，排序融合，RRF，混合检索，候选截断，可复现评价}
\maketitle

\section{阅读问题与证据边界}
\subsection{从一个查询出发}
考虑查询“向量检索中如何融合关键词匹配”。一个词项检索器可能优先返回标题含有“向量”“融合”的文档，另一个语义检索器可能返回没有这些原词、却讨论混合搜索的文章。两者都具有合理依据。把两个列表拼接起来会遇到重复文档；把原始分数相加又默认了分数单位可以比较。排序融合因此不是简单的结果合并，而是在多个不完全一致的相关性证据之间作出决策。

本文关注三个问题。第一，RRF究竟保留了什么、舍弃了什么？第二，当一种方法在经典数据上表现稳定时，我们能把这个经验推广到什么范围？第三，课程项目如何区分真正的融合收益与候选增加、参数搜索或重复检索造成的表面收益？这三个问题把算法公式、实验证据和实现细节连接起来。

\subsection{三篇材料的分工}
主论文为Cormack等人的SIGIR 2009短文\cite{cormack2009rrf}，共两页。它提供RRF的公式、经典实验和简洁的工程论证。补充阅读为Bruch等人的融合函数分析\cite{bruch2023analysis}，本文采用2023年修订的arXiv版本，避免把预印本误写成SIGIR会议论文。第三篇BEIR论文\cite{thakur2021beir}提供跨任务评价的背景，而不提供本文所提项目的实验结果。

本文表格中的手算排序、证明和项目方案均为自主分析。只有明确标为“文献证据”的数值来自原论文。下文的示例文档A、B、C并非真实语料中的检索结果；它们用于检验公式的性质。讨论中出现的潜在收益、失败情形和假设，也不等同于在某个真实数据集上已经验证的结论。

\section{融合前需要明确的检索表示}
\subsection{词项匹配与向量匹配}
TF--IDF可将查询与文档表示为词项权重向量。一种常见定义为
\begin{equation}
 w_{t,d}=\operatorname{tf}(t,d)\log\frac{N}{\operatorname{df}(t)},\qquad
 s_T(q,d)=\frac{\boldsymbol w_q^\top\boldsymbol w_d}
 {\|\boldsymbol w_q\|\|\boldsymbol w_d\|}.
\end{equation}
实际实现往往使用平滑IDF、次线性TF或不同归一化。公式中的选择并非无关紧要：若文档长度处理不同，长文档可能获得不同的排序优势。因此，“TF--IDF基线”应连同分词、IDF和归一化定义一起给出。

BM25同样使用词项证据，但通过词频饱和与长度归一化调整贡献。为统一讨论，本文采用正值IDF的一种常见写法：
\begin{equation}
 s_B(q,d)=\sum_{t\in q}\log\left(1+\frac{N-n_t+0.5}{n_t+0.5}\right)
 \frac{f_{t,d}(k_1+1)}{f_{t,d}+k_1(1-b+b|d|/\overline l)}.
\end{equation}
其中$k_1$控制词频增长的饱和速度，$b$控制长度影响。这里忽略查询词频因子，是为了突出文档端机制；项目若采用其他变体，应记录差异。TF--IDF和BM25并非彼此独立的相关性来源，它们通常共享分词、词项统计和字面匹配偏好。

稠密检索可写为$s_D(q,d)=\boldsymbol e_q^\top\boldsymbol e_d$；当向量单位归一化时等于余弦相似度。它与词项方法的差异在于表示空间及训练目标，而不是名字中带有“向量”就必然更准确。对术语、编号与罕见实体的查询，字面匹配仍可能提供重要证据；反过来，向量模型可能弥补同义表达引起的词汇不匹配。

\subsection{为何原始分数不能直接比较}
假设系统一给A、B的分数分别为9和8，系统二分别给0.2和0.9。相加后A得9.2，B得8.9。现在只把系统二的计分单位乘100：单路次序不变，融合却变为A得29、B得98。直接加分实际上把计分单位当成了系统权重。

这不意味着分数毫无价值。分数间隔可能反映模型对两个候选的区分强度。问题在于必须先说明分数的语义、尺度和归一化，再讨论如何利用它。RRF选择绕开计分单位，代价则是只保留顺序信息。这是一项信息取舍，而不是免费获得的稳健性。

\section{RRF公式及其可证明性质}
\subsection{全排列定义与有限列表实现}
主论文在共同文档集$D$上定义多个排名$r_i$，排名从1开始：
\begin{equation}\label{eq:rrf}
 \RRF_k(d)=\sum_{i=1}^{m}\frac{1}{k+r_i(d)},\qquad k\geq0.
\end{equation}
分数越大，融合名次越靠前。工程上通常只有每路前$L$个文档，不能假装知道列表之外的完整名次。本文约定候选集$C=\bigcup_i C_i$，并定义
\begin{equation}\label{eq:truncated}
 S_k(d)=\sum_{i=1}^{m}\frac{\mathbf1[d\in C_i]}{k+r_i(d)},\quad d\in C.
\end{equation}
这是本文的有限列表实现约定；缺失项贡献0，不能将其误写成原论文全排列定义本身。缺失也不等于已经判定不相关，它只说明该检索器没有在给定预算内返回该文档。

\subsection{单调性、界与不变量}
若某文档在所有系统中的名次都不变或上升，并至少在一路严格上升，则在候选出现情况不变时，其RRF分数严格增加。这直接来自$1/(k+r)$对$r$严格递减。对$m$路有限列表还有
\begin{equation}
 0<S_k(d)\leq\frac{m}{k+1},\quad d\in C.
\end{equation}
这个界说明RRF分数受路数影响，因而跨配置比较原始融合分数通常没有意义。比较的是文档次序与评价指标，不是“0.04比0.03更相关”。

对任意严格递增函数$g_i$，用$g_i(s_i)$替换第$i$路原始分数不会改变单路次序，因而不改变RRF。这解释了它的尺度不变性。它却不具备表示不变性：换分词器、换索引、改变候选截断或引入新的平分规则，都可能改变排名，进而改变融合。

\subsection{头部权重与参数含义}
同一路中第$r$名和第$r+1$名贡献的差为
\begin{equation}
 \Delta_r=\frac{1}{(k+r)(k+r+1)}.
\end{equation}
随着$r$变大，局部位置差异越来越弱；随着$k$增大，所有位置之间的差也变小。第1名相对第$r$名的贡献比为$(k+r)/(k+1)$。当$k=0$时，第1名对第10名的比为10；当$k=60$时则为$70/61$，约1.148。平滑常数因此不是展示层的细节，而是明确改变“单路特别靠前”和“多路共同出现”之间的平衡。

\section{可核验的手算案例与反例}
\subsection{两路互补排序}
构造两路各含三个文档的列表，如表\ref{tab:toy}。设$k=60$，A的得分为$1/61+1/63$，B为$2/62$，C为$1/63$，D为$1/61$。
\begin{table}[ht]
\caption{本文构造的有限列表案例，不是真实实验结果。}
\label{tab:toy}
\centering
\begin{tabular}{ccccc}\toprule
文档 & 路一名次 & 路二名次 & RRF分数 & 融合名次\\\midrule
A&1&3&0.032266&1\\
B&2&2&0.032258&2\\
C&3&缺失&0.015873&4\\
D&缺失&1&0.016393&3\\\bottomrule
\end{tabular}
\end{table}
此例中特别值得解释的是A和B：两者的名次和都是4，RRF却给A略高的分数。令$f(x)=1/(k+x)$，则$f''(x)=2/(k+x)^3>0$，所以$f$为凸函数；相同平均名次下，不平衡的$(1,3)$会比均衡的$(2,2)$获得略高贡献。RRF不是简单平均名次，也不总是偏好“每路都稳定靠前”的候选。

差异只有约$8\times10^{-6}$。若程序先四舍五入到四位小数再排序，会把真实的严格次序变成平分。因此实现时应保留完整浮点值，输出时再格式化，并对真正的同分采用确定性规则，例如文档ID升序。这个细节可以通过单元测试核对，无需真实语料。

\subsection{同一对候选随参数翻转}
构造X的两路名次为$(1,100)$，Y为$(10,10)$。比较
\begin{equation}
 S_k(X)=\frac1{k+1}+\frac1{k+100},\qquad
 S_k(Y)=\frac2{k+10}.
\end{equation}
通分后，$S_k(X)-S_k(Y)$的符号等于$810-81k$的符号。因此$k<10$时X领先，$k=10$时相等，$k>10$时Y领先。这个精确反例不声称真实集合必然如此，却足以否定“改变$k$不会影响排序”的无条件说法。

两份列表中的其余名次可以由不同填充文档占据，故这个构造与完整排名相容。如果只截取前50名，X在第二路不再出现，分数差又发生变化。参数作用与候选深度因而不能分别孤立解释：$k$作用于可见名次，而深度决定哪些证据可见。

\subsection{候选出现频次与大参数极限}
对有限名次范围，当$k$足够大时，可用展开式
\begin{equation}
 \frac1{k+r}=\frac1k-\frac r{k^2}+O(k^{-3})
\end{equation}
得到
\begin{equation}
 S_k(d)=\frac{c(d)}k-\frac{\sum_{i:d\in C_i}r_i(d)}{k^2}
 +O(k^{-3}),
\end{equation}
其中$c(d)$为返回$d$的系统数量。有限列表下，大$k$首先偏好出现次数，然后才区分名次和。若所有系统都给出同一全集的完整排列，所有文档的$c(d)=m$相同，第一项不再区分文档；这时近似比较平均名次。区分这两种输入设定，才能准确解释所谓“共识”。

\subsection{复制检索器为何不是增加知识}
设原来有两路分数贡献$a(d)$与$b(d)$，结果为$a(d)+b(d)$。若把第一路原样复制两次，结果变为$3a(d)+b(d)$。数据与算法都没有新增独立证据，但系统一的相对权重已经变成三倍。

因此，检索器数量不应被等同于信息来源数量。把BM25的几个极相近参数版本全部加入，与加入一个能找到不同相关文档的检索器，机制不同。项目中应把“相同列表重复加入”作为负控制：如果重复使结果改善，能说明特定权重可能有益，却不能说明获得了新互补信息。

\section{与其他融合规则的比较}
\subsection{Borda式排序、Condorcet与CombMNZ}
按名次求和的Borda式方法可以写为$S_{\rm rank}(d)=-\sum_i r_i(d)$，但有限列表的缺失名次需要额外约定。它对头部与尾部的一次位置变化给予同样差值。RRF则用非线性递减函数强化头部，二者不能仅凭“都只用名次”而视为等价。

Condorcet融合通过多系统对候选对的偏好来决定次序。它具有投票式解释，但多数候选对上的偏好不保证存在一个无环的全局顺序。例如三路分别偏好A$>$B$>$C、B$>$C$>$A、C$>$A$>$B，则多数票形成A$>$B、B$>$C、C$>$A的环。具体实现必须处理循环与平票。RRF把各路证据转换为可加分数，避开这一排序关系的不传递性问题，但并未因此取得相关性最优保证。

CombMNZ可概括为返回次数乘以分数和：
\begin{equation}
 S_{\rm MNZ}(d)=c(d)\sum_{i:d\in C_i}\widetilde s_i(d).
\end{equation}
这里用$\widetilde s_i$强调项目需要明确归一化。该规则显式奖励多路命中，同时仍使用分数间距；若计分单位不一致，乘上命中次数不会自动解决尺度问题。对于全为0、常数分数或负分数的输入，还需说明“命中”的判断依据，不能把“分数非零”默认为“确实返回”。

\subsection{凸组合及归一化}
两路分数的凸组合（CC）为
\begin{equation}
 S_{\alpha}(d)=\alpha\widetilde s_1(d)+(1-\alpha)\widetilde s_2(d),
 \quad 0\leq\alpha\leq1.
\end{equation}
常见min--max归一化为
\begin{equation}
 \widetilde s_i(d)=\frac{s_i(d)-s_i^{\min}}
 {s_i^{\max}-s_i^{\min}}.
\end{equation}
如果最大最小值相等，应采用预先规定的退化处理，例如该路对当前查询贡献0，而不是除以0。归一化范围是在全库、候选并集还是单路返回集合上计算，也必须固定。极端高分会压缩其他候选的分数差，所以归一化不仅是代码中的一行变换。

考虑两路正仿射变换$s'_i=a_i s_i+b_i$，其中$a_i>0$且$b_i$对当前查询的所有候选相同。变换后的CC为
\begin{equation}
 S'_{\beta}=\beta a_1s_1+(1-\beta)a_2s_2+\text{常数}.
\end{equation}
若希望其与原来的权重$\alpha$产生相同次序，只需满足
\begin{equation}
 \frac{\beta a_1}{(1-\beta)a_2}=\frac{\alpha}{1-\alpha},\qquad
 \beta=\frac{\alpha a_2}{(1-\alpha)a_1+\alpha a_2}.
\end{equation}
这说明某些尺度变化可以通过权重变换抵消。但如果$a_1,a_2$随查询变化，一个固定的全局$\beta$不一定同时适用于所有查询。另一个限制是候选缺失：对缺失项直接补0，可能破坏“每个候选加同一常数”的条件。讨论归一化等价时必须把假设写清。

\subsection{分数间距确实可能影响合理决策}
自行构造两个候选A、B，系统一始终将A排第一、B排第二。情形甲的系统一归一化分数为$(1,0.99)$，情形乙为$(1,0)$；系统二在两种情形均为$(0,1)$。取$\alpha=0.6$，甲中A得0.6、B得0.994，乙中A得0.6、B得0.4。CC在两种情形给出不同答案，而RRF看到的排名完全相同，不能区分它们。

这个例子不证明CC必然更好，因为分数间距也可能失准。它只说明两种规则可利用的信息不同：RRF不受任意计分尺度的影响，却也无法响应分数中可能存在的强置信度证据。选择规则需要判断这些间距是否可信，以及是否有独立验证查询可以学习权重。

\section{原论文的证据应怎样解读}
\subsection{简短的文献证据摘要}
原论文固定$k=60$用于后续验证；先导表格中$k=80$的MAP为0.2147，$k=60$为0.2145。TREC3中CombMNZ为0.4381，高于RRF的0.4350；TREC9的人工参与最佳系统也高于RRF。LETOR3中RRF的MAP为0.6051，CombMNZ为0.6107，二者差异未达显著（$p=0.2$）\cite{cormack2009rrf}。这些数字是文献结果，并非本文重现实验。

这些证据支持将RRF视为值得保留的简洁基线，却不支持“无需检验就优于任何单路或融合方法”。读论文时，标题、摘要和正文应相互约束：一句范围很宽的总结，必须回到实际比较对象、数据和例外中理解。

\subsection{从经验结论到条件性结论}
“无监督融合”不表示整条检索流水线不需要训练。若输入排名来自监督学习模型，RRF只是在融合阶段无需额外拟合。报告一个免训练融合方法时，应分别描述检索器训练、融合权重学习和测试域适配，避免把不同阶段混为一谈。

“最佳单路”也可能具有两种含义：在独立验证集上选定的部署基线，或事后在测试集合上挑出的最优值。后者适合作为参照上界，但不应作为一个声称无需测试标签就可选择的系统。项目应同时报告每一路，并写清最佳基线如何选择。

单一平均指标可能隐藏失败群体。假设三个查询的基线AP为$(1,0,0)$，融合后为$(0.5,0.3,0.3)$，平均值上升，但原本答对的查询明显退化。真实用户若集中在这类查询，宏平均收益就不足以解释部署价值。因此，除了均值，还应检查逐查询差值、胜负数量和典型失败查询。

\section{现代材料带来的修正}
\subsection{融合函数分析：无需训练不等于无需参数}
Bruch等人在词项与语义检索融合中研究CC与RRF，报告RRF对参数的敏感性，并在其研究设置中观察到调参CC的优势；他们还讨论了线性归一化与可等价排序的权重变换，以及少量标注查询学习融合权重的作用\cite{bruch2023analysis}。这是对特定研究设置的实证与理论分析，不能被改写成“CC在所有任务中都优于RRF”。

对课程项目而言，阅读两篇文献的价值不是给2009年的方法判定过时，而是学习区分三种情境：没有目标域标注时，默认RRF可作为容易启动的方案；有验证数据时，应允许CC和RRF在相同预算下调参；发生域迁移时，源域最优参数是否仍有效需要单独测试。把不同情境合在一个结论里，往往会把数据条件造成的差异归因于公式本身。

理论分析也提示一个报告习惯：指出结论成立的变量范围。式中的秩常数、候选深度与输出深度分别起不同作用，通常都被写成字母$k$会引起混淆。本文将RRF平滑常数记为$k$，单路候选深度记为$L$，最终评价截止位置记为$K$。复现实验的配置文件应沿用这种区分。

\subsection{BEIR：跨域评价的必要性}
BEIR由多个不同领域与任务的数据集组成，原论文使用18个公开数据集评价多种检索架构，显示简单词项基线仍值得重视，效果、计算成本与跨域能力也并不总是同步改善\cite{thakur2021beir}。本文不借用其排行榜来声称课程项目已经具备跨域结论，而是把异构任务视为检查泛化假设的动机。

在一个小型课程项目中，选择一个公开子集完成闭环，比罗列很多没有运行的数据集更可核验。若资源允许，再增加任务差异明显的第二个集合进行外部验证。所有集合的逐查询评价均须使用各自的qrels；不可把某个集合的调参结果与另一个集合的测试标签混合使用。

\section{候选截断、相关性与评价陷阱}
\subsection{任何融合都受候选集上界约束}
设$G_q$为查询$q$的已知相关文档集合，$C_q$为候选并集，$F_q^K\subseteq C_q$为最终前$K$结果，则
\begin{equation}
 \Recall@K(q)=\frac{|F_q^K\cap G_q|}{|G_q|}
 \leq\frac{|C_q\cap G_q|}{|G_q|}.
\end{equation}
右侧是候选集合能达到的召回上限。若相关文档从未被任何一路召回，再精巧的排序公式也无法把它放进结果。评价时应先问是候选发现不足，还是并集内排序失败；两者需要的改进不同。

增大$L$使候选并集单调扩大，但有限列表RRF的最终效果不必单调改善：更多候选会竞争头部位置，既有候选也可能在另一路新获得一项贡献。公平实验应分别比较“同每路深度”和“同总候选预算”两种设置，说明融合是否额外用了更多检索计算。

\subsection{重叠不等于互补，也不等于无用}
可用Jaccard系数描述返回集合的重叠：
\begin{equation}
 J(C_1,C_2)=\frac{|C_1\cap C_2|}{|C_1\cup C_2|}.
\end{equation}
低重叠并不保证收益，因为两路可能只是各自找到不同的不相关文档。真正有用的互补可以通过独有相关文档数考察，例如$|(C_2\setminus C_1)\cap G_q|$。高重叠也不意味着一路多余，因为两路可能在同一候选集内提供不同次序，有助于重排。

次序相关性可以在共同候选上计算，但必须说明缺失项处理方式。只在交集上计算相关性，会忽略恰好最能体现召回差异的部分；给所有缺失项同一个末位，也会制造大量平分。项目中将集合重叠、独有相关文档与最终效果分开报告，比把它们压成一个“多样性分数”更容易解释。

\subsection{指标的含义与边界}
对二值相关性，平均精确率可写为
\begin{equation}
 \AP(q)=\frac1{|G_q|}\sum_{j=1}^{n}P@j(q)\,\Rel_q(j),
 \qquad \operatorname{MAP}=\frac1{|Q|}\sum_q\AP(q).
\end{equation}
如果只保存前$n$个结果，未返回的相关项贡献0。若按cutoff报告AP，应把该cutoff明确写出。MAP对所有查询等权，不能与把查询的相关文档混在一起计算的微平均混用。

分级相关性常用
\begin{equation}
 \DCG@K(q)=\sum_{j=1}^{K}\frac{2^{g_q(j)}-1}{\log_2(j+1)},
 \quad \operatorname{nDCG}@K=\frac{\DCG@K}{\operatorname{IDCG}@K}.
\end{equation}
应预先规定无相关项查询或IDCG为0时的处理。MRR则只看首个相关结果的位置倒数；它不是RRF的缩写变体，也不能反映多个相关文档的完整覆盖。本文建议头部nDCG、Recall和MAP共同报告，避免一个指标包办所有质量维度。

\subsection{不完全标注与统计解释}
qrels是已获得的相关性判断，而非完备真值。通常把未标注文档按不相关计分，便于统一计算，但新的召回器找到未判断文档时可能被低估。课程报告应保留这个限制，不擅自修改官方qrels来让某种方法显得更好。个例分析可另外记录“未判断”，与已确认不相关区分。

对两个系统，令$\delta_q=M_q(A)-M_q(B)$，先报告平均差值和查询级分布，再考虑统计检验。一个透明的方案是按查询配对重采样bootstrap，报告均值差的置信区间、重复次数和随机种子。若同时搜索很多参数，测试集上最好的差值会有选择偏差；因此参数选择在验证集完成，测试集仅做一次最终比较。置信区间跨0不表示两个系统完全相同，只表示现有样本不足以排除零差异。

\section{课程项目：把阅读问题变为可执行协议}
\subsection{研究假设与最小闭环}
项目可命名为“公开语料上的多路排序融合及敏感性分析”。最小版本保留TF--IDF、BM25两路，比较单路、RRF和归一化CC；可选加入一个公开稠密检索器形成更有差异的来源。若没有运行稠密模型，不应把词项两路的结果包装成已经完成的语义混合检索。

拟检验的假设为：RRF能否在固定检索预算下改善头部排序；收益是否与独有相关候选有关；默认常数60与验证集选择的常数是否有稳定差异；归一化CC是否通过少量验证查询获得收益。这些均为待检验假设，而不是报告预设结论。负结果同样可以回答问题，例如词项两路过于相似，融合没有改善，就是有解释价值的发现。

\subsection{数据、划分与泄漏防护}
输入由corpus、queries和qrels三部分组成，采用稳定文档ID与查询ID连接。语料预处理应确定Unicode规范化、大小写、标点、停用词与词干化的策略；并保留文档标题是否拼接正文的信息。所有检索器使用同一可检索文档集合，避免某一路少建了文档而造成伪差异。

若集合提供官方训练、验证、测试划分，优先遵循。若没有，则在查看测试表现前按固定种子划分有标注查询，并保存ID列表。小数据上可采用验证折选择参数，保留一个最终测试部分。不能用测试集最优$k$宣称“默认RRF效果”，也不能先挑选融合改善最多的查询再展示整体结论。

\subsection{参数与对照实验}
表\ref{tab:protocol}给出计划中的参数轴。网格规模应受运行预算限制，所有候选配置与选择规则提前保存。参数之间可能有相互作用，首先固定其他参数做一维曲线，再在独立验证集选择最终配置，减少解释上的自由度。
\begin{table}[ht]
\caption{建议实验协议；本表仅为计划，不包含完成结果。}
\label{tab:protocol}
\centering\small
\begin{tabular}{p{0.22\columnwidth}p{0.68\columnwidth}}\toprule
轴或对照&预先规定的检查方式\\\midrule
单路基线&TF--IDF、BM25；可选稠密检索\\
RRF常数&$k\in\{0,10,20,60,100\}$；60单独列为默认\\
候选深度&$L\in\{10,50,100\}$，输出$K$固定\\
CC权重&$\alpha\in\{0,0.1,\ldots,1\}$，仅验证集选择\\
控制变量&同语料、同查询、同qrels、同分词版本\\
负控制&复制一路列表；改变输出格式不改变分数\\
结果呈现&nDCG@10、Recall@100、MAP与逐查询差值\\\bottomrule
\end{tabular}
\end{table}

同一参数网格在不同数据集上是否需要重新调参，应作为实验设置的一部分。如果为了目标域性能重新选择参数，就不能再把该结果称为完全零样本。端点$\alpha=0,1$也应保留，防止融合网格排除了单路而迫使系统“必须融合”。

\subsection{实现、核对与运行产物}
一个清晰的实现流程为：读入语料并建立索引；每路对每个查询输出带docid、rank、score的运行文件；对相同query连接候选；计算融合；用同一评价器运行；导出逐查询与汇总表。将检索和融合分离后，参数实验可以复用已有运行文件，不必每次重建索引。

RRF逐路遍历返回项，用字典保存文档ID与累加分数。累加的时间复杂度为$O(mL)$。若对候选并集完全排序，整体时间成本为
\begin{equation}
 T_{\mathrm{fusion}}=O\bigl(mL+|C|\log|C|\bigr).
\end{equation}
这些仅是融合成本，不含检索、编码或索引构建。报告应区分冷启动、索引构建和查询阶段。把缓存后的查询时间与全流程时间混比，会夸大效率差异。

最小正确性核对应包括：排名从1开始；同一路重复docid不重复计票；缺失不贡献；同分次序可复现；手算表\ref{tab:toy}匹配；复制一路等价于增加对应权重；在所有路相同且文档集合一致时，融合次序不变。评价器另需用一个人工标注的小集合检查AP与nDCG，避免把实现错误当成算法结论。

应保存的产物包括配置、环境版本、输入下载来源、数据校验信息、固定划分、单路run文件、融合run文件、逐查询指标、汇总结果和运行日志。结果表中的空白只能解释为未运行，不能填入原论文数值作为自己的复现分数。若加入稠密模型，还需记录模型名称与版本、向量归一化、相似度和近邻搜索是否精确。

\subsection{结果解释的优先顺序}
完成实验后先检查单路与候选召回上界，再比较融合质量，随后解释参数和失败查询。若融合未提升，应先排除索引、ID对齐和评价错误，再检查两路相关性过高或某一路质量太差。若融合提升，则检查是否只因候选增多，是否在固定预算下仍成立，是否对少数查询损伤严重。

展示时可选一个改善查询与一个退化查询，列出原始名次、融合贡献、相关性标签及文档片段。例子应来自预先明确的选择标准，例如最大正差和最大负差，避免只挑支持观点的案例。最终结论应写成“在指定数据、预算和指标下观察到什么”，而不是把某个默认参数升级为通用规律。

\section{进一步分析：从排序行为到决策代价}
\subsection{文档级加分不能替代查询级可靠性判断}
RRF对所有查询采用相同的计分函数。一条查询可能是精确编号，另一条可能是概念描述；两类查询对词项与语义证据的依赖可能不同。如果对某类查询，系统一的排名更可靠，等权融合就存在把可靠结果拉向不可靠结果的风险。这个问题不能只靠增加一个全局平滑常数解决，因为常数同时作用于所有候选和查询。

一个可分析的扩展是加权RRF：
\begin{equation}
 S_{k,w}(q,d)=\sum_i w_i(q)\frac{\mathbf1[d\in C_i(q)]}{k+r_i(q,d)},
 \quad w_i(q)\geq0.
\end{equation}
若所有权重相同，则退化为等权版本；将权重统一乘正常数不会改变排序，因此可令$\sum_iw_i(q)=1$。需要强调的是，加权版本并非原始RRF不调参基线：如果$w_i(q)$依赖标签学习，就应计算训练样本与特征获取成本，并与一个只学习全局权重的简单对照比较。

查询自适应并不应成为课程项目的必要任务。先验证全局等权、全局可学习权重的差异，已经能够回答“性能限制来自证据本身还是等权假设”的一部分。若直接引入复杂门控网络，参数、训练与检索表示的变化会叠在一起，难以归因。

\subsection{弱检索器能否被自动忽略}
设原有系统对A、B的融合分差为$\gamma=S(A)-S(B)>0$，新增一路的贡献差为
\begin{equation}
 \eta=\frac{\mathbf1[A\in C_{m+1}]}{k+r_{m+1}(A)}
 -\frac{\mathbf1[B\in C_{m+1}]}{k+r_{m+1}(B)}.
\end{equation}
加入后若$\gamma+\eta<0$，A、B的次序就翻转。由于$|\eta|\leq1/(k+1)$，可以得到一个充分条件：若$\gamma>1/(k+1)$，单路新增证据无法反转这对候选。然而多数真实头部候选的差距可能远小于这个界，是否安全不能仅凭系统数量判断。

若新增一路只给B高排名而完全漏掉A，最大负扰动就接近$-1/(k+1)$。RRF不知道这一路很弱；“它只用名次”也意味着“它无法直接利用分数是否经过校准来识别弱源”。这一分析提示应同时评价每路单独效果、移除一路的消融与加入噪声路的压力测试。只展示最佳融合组合，无法判断方法对低质量来源是否稳健。

\subsection{并集、交集与推荐解释不同}
对两个有限候选列表，选择并集允许任一路独有的文档进入最终结果，选择交集则要求双方共同返回。交集可以减少候选数量，却会把互补召回的潜在价值直接删除。若相关文档只出现在一路，交集上的任何排序都无法保留它。因此两种候选策略对应不同决策，不能在比较两种融合公式时偷偷同时更改候选策略。

另一方面，多路共同命中只能提供一种一致性证据，并不等于独立验证。一份错误文档若因共同分词或热门词而被所有词项检索器返回，RRF仍可能给它很高分。系统对外展示“被三种检索器支持”时，应避免把算法数量包装成三个独立可信来源。若还要生成答案，文档事实的可核验性属于另一层问题，排序分数本身并不能完成事实验证。

\subsection{平分排名与确定性处理}
单路中经常出现相同分数，尤其是短查询、短文本或离散计分的情况。可以给平分文档分配平均秩，也可以按稳定文档ID打破平分。两者对RRF不完全等价，因为倒数函数为非线性。若A、B共享位置1与2，平均秩方案给两者$1/(k+1.5)$；稳定ID方案则分别给$1/(k+1)$与$1/(k+2)$。

平均秩表达“证据不足以区分”，稳定规则则便于生成严格列表。课程项目可选择其中一种，但应在所有算法中一致处理，且保存该选择。如果只有输出文件的存储顺序改变，就导致最终评价改变，通常意味着某处把非确定遍历顺序当成了排序规则。稳定结果不仅方便复现，也能避免把微小随机波动误解释为参数差异。

\section{一个完整的人工评价演练}
\subsection{从名次变化计算指标变化}
为检查评价定义，另构造一个与真实语料无关的查询。设其相关文档恰好为A和C，单路前四项依次为B、A、D、C，融合结果为A、B、C、D。所有文档均已有判断，且只有A、C相关。单路的AP为
\begin{equation}
 \AP_{\rm single}=\frac12\left(\frac12+\frac24\right)=0.5,
\end{equation}
而融合AP为
\begin{equation}
 \AP_{\rm fused}=\frac12\left(1+\frac23\right)=\frac56.
\end{equation}
Recall@4在两者中均为1；Recall@2也均为$1/2$。这些召回指标不变，不代表排序没有改善，而是它们只检查在截止位置前覆盖了多少相关文档，不衡量相关文档在窗口内的具体次序。

采用二值gain并取$K=3$，单路的DCG为$1/\log_2 3$，融合为$1+1/\log_2 4=1.5$。理想排序为A、C、B，因此
\begin{align}
 \operatorname{IDCG}@3&=1+\frac1{\log_2 3},\\
 \operatorname{nDCG}_{\rm single}@3&\approx0.3869,\\
 \operatorname{nDCG}_{\rm fused}@3&\approx0.9197.
\end{align}
这些数值是由人工构造名单直接计算的演练，不是本文项目在公开数据上的成绩。它们的作用是提供一个可手动审计的评价器输入与预期输出。

\subsection{为什么只报告最终得分不够}
如果只看到MAP从0.5变成$5/6$，无法知道改善来自检索了新文档还是重新排列了原候选。在此演练中，候选集完全一致，只改变次序，所以可把收益归于重排。真实项目必须同时保存候选与排序，才能作出同样清楚的归因。

进一步，如果C在官方qrels中未被判断而按0计，则该演练的可观察相关集合发生变化，指标也会不同。算法可能未变，评价依据却变了。因而运行结果文件和qrels版本应共同构成一个评价单元；共享某个分数但不共享判断文件与指标定义，无法构成完整复现。

\subsection{聚合后的平均值可能如何误导}
设两个查询的nDCG变化分别为$+0.4$与$-0.3$，宏平均提升0.05。若第一个查询每天只出现一次，第二个查询每天出现一百次，真实访问加权变化则约为$-0.293$。课程数据集通常采用查询等权，是合理且可重复的基准约定；但它与生产访问分布不是同一个目标。

这不要求在课程报告中虚构真实访问量，而是要求写清评价对象。可以说“在测试查询等权条件下提高平均质量”，不宜直接说“所有用户体验改善”。若将来有真实访问日志，还需考虑隐私、查询重复、时效与流量偏差，另行设计权重和评估流程。

\subsection{预注册式记录能减少解释自由度}
在运行前写下主要指标、基线、参数网格、查询划分与失败处理，是一种低成本的自我约束。例如主要指标定为nDCG@10，Recall@100和MAP作为辅助；选择验证集主要指标最高的配置，若相同则优先更小候选深度，再优先更简单参数。这样可以避免在看到结果后临时更换指标，让一组偶然改善成为主要发现。

报告仍可以展示探索性发现，但应把预先计划与事后分析分开。若某个查询类别在结果出现后才被识别出来，就应称为探索性案例，而不是已经独立证实的泛化规律。严格限制表述范围不会削弱项目，反而使读者能够准确判断证据的强弱。

\section{结语：如何形成可辩护的阅读结论}
RRF值得学习，不仅因为公式简短，更因为它展示了一种明确的工程取舍：用序位证据换取对任意计分尺度的独立性。本文的自主构造说明，参数能够改变排名，候选截断能够改变证据，重复来源能够改变隐含权重，而忽略分数也会忽略有价值的间距信息。这些性质共同决定了它应被当作需要评价的基线，而不是默认正确的答案。

对课程学习而言，最有意义的工作是把可解释公式、反例和可复现评价连接起来。三篇材料分别帮助提出算法、限制推广范围和设计外部检验；项目则需要用独立的真实运行证据完成最后一步。本文已完成的是阅读与分析，关于具体语料上的效果、延迟和最优参数，须由项目实测后再作判断。

\paragraph{AI辅助说明}
本文使用AI辅助检索文献、推导检查与组织表达。参考文献和引用数据均标明来源；构造案例不作为真实实验。提交者应核对全文、理解公式和局限，并据实际完成的实验补充说明后署名。

\bibliographystyle{ACM-Reference-Format}
\bibliography{references}
\end{document}
